% Makro bantu. Dimuat setelah packages.tex, colors.tex, facts.tex.

% --- Metadata dokumen ------------------------------------------------------
\newcommand{\reporttitle}{Laporan Desain Sistem Pemberi Pakan Ikan Otomatis}
\newcommand{\reportauthor}{Proyek Akhir — Autofeeder Aquaponik}
\newcommand{\reportsubject}{ESP32, MQTT, SQLite, Streamlit}

\hypersetup{
  pdftitle={\reporttitle},
  pdfauthor={\reportauthor},
  pdfsubject={\reportsubject},
  pdfkeywords={ESP32, MQTT, SQLite, Streamlit, autofeeder, aquaponik}
}

% --- Mode keluaran ---------------------------------------------------------
% Diatur oleh wrapper: \standalonetrue (mandiri) atau \standalonefalse (lampiran).
% (dideklarasikan di wrapper sebelum preamble dimuat)

% --- Teks teknis -----------------------------------------------------------
% \code memakai \detokenize sehingga garis bawah dan kurung tidak perlu di-escape:
%   \code{feed_level_pct}  -> feed_level_pct
\newcommand{\code}[1]{\texttt{\detokenize{#1}}}

% Topik MQTT lengkap: \topic{telemetry} -> feeder-ayurahmaw/feeder-64008c/telemetry
\newcommand{\topic}[1]{\texttt{\TopicPrefix/\DeviceExample/\detokenize{#1}}}

% Topik dengan segmen bebas: \topicof{feeder01/ack}
\newcommand{\topicof}[1]{\texttt{\TopicPrefix/\detokenize{#1}}}

% Nama berkas / path.
\newcommand{\file}[1]{\texttt{\detokenize{#1}}}

% --- Gambar ---------------------------------------------------------------
% Hasil render Mermaid (SVG -> PDF) ada di assets/diagrams/pdf/
\newcommand{\mmdfig}[1]{assets/diagrams/pdf/#1}
% Hasil compile TikZ standalone ada di assets/tikz/build/
\newcommand{\tikzfig}[1]{assets/tikz/build/#1}

% Gambar standar: \figureasset[lebar]{berkas}{caption}{label}
% Tinggi dibatasi 0.88\textheight agar diagram yang tinggi tidak meluber ke halaman
% berikutnya; keepaspectratio menjaga proporsi.
\newcommand{\figureasset}[4][\linewidth]{%
  \begin{figure}[H]
    \centering
    \includegraphics[width=#1,height=0.88\textheight,keepaspectratio]{\mmdfig{#2}}
    \caption{#3}
    \label{#4}
  \end{figure}%
}
\newcommand{\figuretikz}[4][\linewidth]{%
  \begin{figure}[H]
    \centering
    \includegraphics[width=#1,height=0.88\textheight,keepaspectratio]{\tikzfig{#2}}
    \caption{#3}
    \label{#4}
  \end{figure}%
}

% Rujukan singkat.
\newcommand{\figref}[1]{Gambar~\ref{#1}}
\newcommand{\tabref}[1]{Tabel~\ref{#1}}
\newcommand{\chapref}[1]{Bab~\ref{#1}}

% --- Provenance -----------------------------------------------------------
% Jejak ke kode: \src{firmware/esp32\_feeder/config.h}  (pakai \detokenize)
\newcommand{\src}[1]{\par\smallskip{\footnotesize\color{muted}Sumber: \texttt{\detokenize{#1}}}\par}
